\section*{Bab \ref{ch:b1:flowering-plant-organization} --- Organisasi Fungsional Sebuah Tumbuhan Berbunga}

\begin{solution}{exo:b1:flowering-plant-organization:1}
\omterm{def:b1:flowering-plant-organization:organs}{Akar}: penambatan, penyerapan air dan ion mineral. \omterm{def:b1:flowering-plant-organization:organs}{Batang}: penopang \omterm{def:b1:flowering-plant-organization:organs}{daun}
ke arah cahaya, penghantaran antara \omterm{def:b1:flowering-plant-organization:organs}{akar} dan \omterm{def:b1:flowering-plant-organization:organs}{daun}. \omterm{def:b1:flowering-plant-organization:organs}{Daun}: fotosintesis
(dan pertukaran gas, transpirasi).
\end{solution}

\begin{solution}{exo:b1:flowering-plant-organization:2}
Satu modul adalah satu \omterm{def:b1:flowering-plant-organization:organs}{buku} dengan \omterm{def:b1:flowering-plant-organization:organs}{daun} dan tunas ketiaknya, ditambah
\omterm{def:b1:flowering-plant-organization:organs}{ruas} di bawahnya. \omterm{prop:b1:flowering-plant-organization:modular}{Pertumbuhan tak tentu}: tidak ada ukuran dewasa; tunas
ujung dan tunas ketiak terus menambah modul selama tumbuhan itu hidup.
\end{solution}

\begin{solution}{exo:b1:flowering-plant-organization:3}
\omterm{def:b1:flowering-plant-organization:tissues}{Epidermis}: hidup; penyelimut, \omterm{def:b1:flowering-plant-organization:tissues}{kutikula}, \omterm{def:b1:flowering-plant-organization:tissues}{stomata}, \omterm{def:b1:flowering-plant-organization:organs}{rambut akar}. \omterm{def:b1:flowering-plant-organization:tissues}{Parenkim}:
hidup; fotosintesis, penimbunan. \omterm{def:b1:flowering-plant-organization:tissues}{Kolenkim}: hidup; penopang lentur.
\omterm{def:b1:flowering-plant-organization:tissues}{Sklerenkim}: mati; penopang kaku. Sel penghantar \omterm{def:b1:flowering-plant-organization:tissues}{xilem}: mati; pengangkutan
air. \omterm{def:b1:flowering-plant-organization:tissues}{Pembuluh tapis} \omterm{def:b1:flowering-plant-organization:tissues}{floem}: hidup (dengan sel pengiringnya); pengangkutan
gula.
\end{solution}

\begin{solution}{exo:b1:flowering-plant-organization:4}
$e^{-0.9} = 0.41$, yakni \qty{41}{\%}; $e^{-3} = 0.05$, yakni
\qty{5}{\%}.
\end{solution}

\begin{solution}{exo:b1:flowering-plant-organization:5}
Sebuah \omterm{def:b1:flowering-plant-organization:organs}{akar} (satu silinder pusat, \omterm{prop:b1:flowering-plant-organization:stemroot}{endodermis}, \omterm{def:b1:flowering-plant-organization:tissues}{xilem} dan \omterm{def:b1:flowering-plant-organization:tissues}{floem}
berselang-seling pada satu jari-jari) tumbuhan dikotil (sedikit lengan
\omterm{def:b1:flowering-plant-organization:tissues}{xilem}, tanpa empulur di pusat).
\end{solution}

\begin{solution}{exo:b1:flowering-plant-organization:6}
Sebuah balok yang dilenturkan mengalami tegangan terbesar di
permukaannya dan tidak sama sekali sepanjang sumbunya, sehingga bahan
yang ditempatkan pada cincin tepi menahan lenturan dengan massa paling
kecil (sebuah pipa). Sebuah kabel yang ditarik mengalami tegangan merata
pada penampangnya dan melentur bebas; untai di pusat memikul tarikannya
sementara korteks tetap lentur, dan \omterm{def:b1:flowering-plant-organization:organs}{akar} samping dapat menembusnya.
\end{solution}

\begin{solution}{exo:b1:flowering-plant-organization:7}
$1 - e^{-2.4} = 0.91$; pada $L = 6$, $1 - e^{-3.6} = 0.97$: dua lapis
\omterm{def:b1:flowering-plant-organization:organs}{daun} lagi menambah \qty{6}{\%} cahaya sementara memakan biaya
\qty{50}{\%} \omterm{def:b1:flowering-plant-organization:organs}{daun} lebih banyak.
\end{solution}

\begin{solution}{exo:b1:flowering-plant-organization:8}
$S = \pi d L$: $d = 240/(\pi\times 623\,000) = \qty{1.2e-4}{m}$, sekitar
\qty{0.12}{mm}. $(240 + 400)/5 = 128$.
\end{solution}

\begin{solution}{exo:b1:flowering-plant-organization:9}
Permukaan \omterm{def:b1:flowering-plant-organization:organs}{akar} tiap mm: $\pi\times 0.4 = \qty{1.26}{mm^2}$. Rambut tiap
mm: $250\times\pi\times 0.010\times 0.8 = \qty{6.3}{mm^2}$. Seluruhnya
$7.5/1.26 = 6$: rambutnya melipatgandakan permukaan enam kali.
\end{solution}

\begin{solution}{exo:b1:flowering-plant-organization:10}
Di dalam hutan tunas ketiak yang rendah ternaungi; modul yang akan
dibuatnya tidak dapat membayar dirinya sendiri, sehingga tunas itu
tinggal tidur atau cabangnya mati, dan pertumbuhannya dialokasikan ke
tunas ujung yang berlomba naik ke arah cahaya. Di ladang tiap tunas
terkena cahaya dan tiap cabang membayar, sehingga tajuknya terisi penuh
sampai bawah. Biayanya: \omterm{def:b1:flowering-plant-organization:organs}{batang} jangkung yang kurus adalah tanaman modal
besar dalam kayu yang tak menghasilkan, rentan terhadap angin;
keuntungannya: tumbuhan itu menyesuaikan bentuknya dengan tempat cahaya
berada, tanpa rencana yang tetap.
\end{solution}

\begin{solution}{exo:b1:flowering-plant-organization:11}
Lekuk dan rambutnya menahan selapis udara lembap yang diam di atas
liangnya; gradien uap dari ruang udara ke udara yang bergerak terbentang
sepanjang lintasan yang lebih panjang dan transpirasinya turun. Karbon
dioksida harus berdifusi masuk sepanjang lintasan panjang yang sama,
sehingga penyerapannya juga turun, meskipun lebih sedikit, karena
gradien $\mathrm{CO_2}$ ditentukan atmosfer dan pemakaian \omterm{def:b1:flowering-plant-organization:organs}{daunnya}. Di
gurun air adalah sumber daya yang membatasi dan pertukaran itu
menyelamatkan hidup tumbuhannya; di hutan hujan air gratis dan perolehan
karbonlah yang diganjar persaingan.
\end{solution}

\begin{solution}{exo:b1:flowering-plant-organization:12}
Tumbuhan tidak memiliki \omterm{def:b1:mammal-organization:compartments}{cairan ekstrasel} teratur yang dijaga pada suhu
dan susunan tetap antara selnya dan dunia luar: selnya menemui air
tanah, udara dan sinar matahari secara langsung, dan suhunya adalah suhu
udara. Karena itu tiap selnya mengatur isinya sendiri (ion, air, pH)
melintasi membran dan dindingnya; tumbuhan itu mengatur pertukarannya
pada permukaan (\omterm{def:b1:flowering-plant-organization:tissues}{stomata} membuka dan menutup, \omterm{prop:b1:flowering-plant-organization:stemroot}{endodermis} memilih apa yang
masuk ke \omterm{def:b1:flowering-plant-organization:tissues}{xilem}) dan menyetel tubuhnya lewat pertumbuhan (lebih banyak
\omterm{def:b1:flowering-plant-organization:organs}{akar} di tanah kering, lebih banyak \omterm{def:b1:flowering-plant-organization:organs}{daun} di tempat teduh). Ia memang
mengatur, tetapi pengaturannya ada pada permukaan dan pada bentuk, bukan
di dalam sebuah cairan pribadi.
\end{solution}

\begin{solution}{pb:b1:flowering-plant-organization:1}
\textbf{1.} $5\times 100 = \qty{500}{m^2}$.
\textbf{2.} $e^{-2.5} = 0.082$ mencapai tanah; \qty{92}{\%} ditangkap.
\textbf{3.} $L = 2$: $1 - e^{-1} = \qty{63}{\%}$; $L = 8$: $1 -
e^{-4} = \qty{98}{\%}$.
\textbf{4.} Dari 2 ke 5 tajuknya memperoleh \qty{29}{\%} cahaya untuk
\qty{300}{m^2} \omterm{def:b1:flowering-plant-organization:organs}{daun}; dari 5 ke 8 ia hanya akan memperoleh \qty{6}{\%}
untuk \qty{300}{m^2} lagi yang biaya pembangunan dan pemasokan airnya
sama besar: \omterm{def:b1:flowering-plant-organization:organs}{daun} tambahan itu tidak membayar dirinya.
\textbf{5.} $500/0.005 = \num{100000}$ helai \omterm{def:b1:flowering-plant-organization:organs}{daun}.
\textbf{6.} $500\times 10^6\,\mathrm{mm^2}\times 150 = \num{7.5e10}$
\omterm{def:b1:flowering-plant-organization:tissues}{stoma}.
\textbf{7.} $\qty{50}{\micro m^2}$; tiap mm$^2$: $150\times 50\times
10^{-6} = \qty{7.5e-3}{mm^2}$, yakni \qty{0.75}{\%} permukaannya.
\textbf{8.} $500\times 4\times 10^{-6} = \qty{2e-3}{mol/s}$, yakni
\qty{88}{mg/s}.
\textbf{9.} $0.088\times 43\,200 = \qty{3.8}{kg}$ $\mathrm{CO_2}$, yang
mengandung $3.8\times 12/44 = \qty{1.04}{kg}$ karbon.
\textbf{10.} $150\times 1.04 = \qty{156}{kg}$ karbon; separuhnya,
\qty{78}{kg}, menjadi kayu berkadar karbon \qty{50}{\%}: \qty{156}{kg}
kayu.
\textbf{11.} $500\times 10^{-3} = \qty{0.5}{mol/s}$, yakni \qty{9}{g/s}.
\textbf{12.} $9\times 43\,200 = \qty{389}{kg}$, sekitar \qty{390}{L}.
\textbf{13.} $0.5/(2\times 10^{-3}) = 250$ molekul air tiap
$\mathrm{CO_2}$. Liang yang meloloskan $\mathrm{CO_2}$ masuk
meloloskan air keluar; udara di dalam jenuh sedangkan $\mathrm{CO_2}$
hanya \qty{0.04}{\%} udara luar, sehingga gradien uap ke luar jauh lebih
curam daripada gradien $\mathrm{CO_2}$ ke dalam.
\textbf{14.} $\qty{2}{mm}\times\qty{100}{m^2} = \qty{200}{L}$: setengah
hari.
\textbf{15.} \qty{100}{m^3}; $5\times 100 = \qty{500}{km}$.
\textbf{16.} $\pi\times 0.5\times 10^{-3}\times 5\times 10^5 =
\qty{785}{m^2}$.
\textbf{17.} $200\times 5\times 10^8\,\mathrm{mm} = \num{1e11}$ rambut.
\textbf{18.} Masing-masing $\pi\times 10^{-5}\times 5\times 10^{-4} =
\qty{1.57e-8}{m^2}$; seluruhnya \qty{1570}{m^2}.
\textbf{19.} $785 + 1570 = \qty{2355}{m^2}$, $4.7$ kali luas \omterm{def:b1:flowering-plant-organization:organs}{daunnya}.
\textbf{20.} Volume dalam \qty{1}{mm}: $\pi\times(10^{-3})^2\times 5
\times 10^5 = \qty{1.6}{m^3}$, yakni \qty{1.6}{\%} tanahnya.
\textbf{21.} Gandum hitam: $\pi\times 10^{-6}\times 6.2\times 10^5 =
\qty{1.95}{m^3}$ dari kotak \qty{0.05}{m^3}, yakni \qty{3900}{\%}: tiap
titik tanahnya berada dalam satu milimeter dari empat puluh \omterm{def:b1:flowering-plant-organization:organs}{akar}.
Gandum hitam itu menjelajahi tanahnya lebih dari dua ribu kali lebih
saksama --- tanaman pertanian dengan anggaran empat bulan berhadapan
dengan sebatang pohon yang menjelajahi seratus meter kubik.
\textbf{22.} $1000 + 2355 = \qty{3355}{m^2}$; $\qty{34}{m^2}$ tiap
meter persegi tanah.
\textbf{23.} Manusia: $130/2 = 65$; rasio pohon ek itu separuh rasio
permukaan terlipat manusia, tanpa pelipatan sama sekali.
\textbf{24.} Pohon ek menambah modul, \omterm{def:b1:flowering-plant-organization:organs}{daun} dan \omterm{def:b1:flowering-plant-organization:organs}{akar} tiap tahun, sehingga
permukaannya tumbuh bersama umurnya; permukaan mamalia adalah \omterm{def:b1:mammal-organization:organ}{organ}
yang tetap, lengkap pada ukuran dewasa, dan satu-satunya cara menaikkan
rasionya adalah melipatnya selama perkembangan.
\textbf{25.} Sekitar \qty{34}{m^2} \omterm{prop:b1:organism-environment:surfaces}{permukaan pertukaran} tiap meter
persegi tanah: \qty{10}{m^2} \omterm{def:b1:flowering-plant-organization:organs}{daun} (kedua sisinya) di atas dan
\qty{24}{m^2} \omterm{def:b1:flowering-plant-organization:organs}{akar} dan \omterm{def:b1:flowering-plant-organization:organs}{rambut akar} di bawah.
\end{solution}
